\begin{theorem} (б/д)
    В условиях (R0)-(R8) любая состоятельная последовательность оценок $\displaystyle \hat{\theta }_{n}( X_{1} ,\ \dotsc ,\ X_{n})$, являющихся решениями уравнения правдоподобия, удовлетворяют соотношению
    \begin{equation*}
        \forall \theta \in \Theta \hookrightarrow \sqrt{n}(\hat{\theta }_{n} -\theta )\xrightarrow{d_{\theta }}\mathcal{N}\left( 0,\, \dfrac{1}{i( \theta )}\right) ,
    \end{equation*}
    где $\displaystyle i( \theta )$ -- количество информации Фишера в одном наблюдении.
\end{theorem}

\begin{note}
    Формально перенеся $\displaystyle \sqrt{n}$ в правую часть, получим, что асимптотическая дисперсия будет равняться информации Фишера всей выборки. Таким образом, эту теорему можно воспринимать, как асимптотическую версию неравенства Рао-Крамера.
\end{note}
\begin{corollary}
    (асимптотическая нормальность ОМП) В условиях (R0)-(R8), если $\displaystyle \forall n\in \mathbb{N} \ \forall X_{1} ,\ \dotsc ,\ X_{n}$ существует единственное решение $\displaystyle \hat{\theta }_{n}$ уравнения правдоподобия, и если оно является статистикой (т.е. измеримой функцией выборки), то $\displaystyle \hat{\theta }_{n}$ -- асимптотически нормальная оценка $\displaystyle \theta $ с асимптотической дисперсией $\displaystyle \dfrac{1}{i( \theta )}$ и с вероятностью, стремящейся к единице, является ОМП. 
\end{corollary}
\begin{theorem}
    (Бахадур, б/д) Пусть выполнены условия регулярности (R0)-(R8), и $\displaystyle \hat{\theta }_{n}( X_{1} ,\ \dotsc ,\ X_{n})$ асимптотически нормальная оценка $\displaystyle \theta $, т.е. $\displaystyle \sqrt{n}(\hat{\theta }_{n} -\theta )\xrightarrow{d_{\theta }}\mathcal{N}\left( 0,\ \sigma ^{2}( \theta )\right)$, причем $\displaystyle \sigma ( \theta )$ непрерывна. Тогда, $\displaystyle \forall \theta \in \Theta \hookrightarrow \sigma ^{2}( \theta ) \geqslant \dfrac{1}{i( \theta )}$.
\end{theorem}
\begin{note}
    Эта теорема также является некоторой асимптотической версией неравенства Крамера-Рао.
\end{note}
\begin{corollary}
    В условиях следствия про асимптотическую нормальность ОМП имеем, что ОМП не хуже любой другой оценки в асимптотическом подходе в классе оценок с непрерывной дисперсией.
\end{corollary}
\begin{definition}
    Если $\displaystyle \sqrt{n}(\hat{\theta }_{n}( X_{1} ,\ \dotsc ,\ X_{n}) -\theta )\xrightarrow{d_{\theta }}\mathcal{N}\left( 0,\ \dfrac{1}{i( \theta )}\right)$, то $\displaystyle \hat{\theta }_{n}$ называется асимптотически эффективной оценкой $\displaystyle \theta $.
\end{definition}
\begin{proposition}
    Пусть выполнены условия регулярности для неравенства Крамера-Рао, $\displaystyle \hat{\theta }( X)$ -- эффективная оценка $\displaystyle \theta $, и равенство из критерия эффективности (для $\displaystyle \hat{\theta }( X)$) выполнено для любого $\displaystyle X$ и любого $\displaystyle \theta $. Тогда $\displaystyle \hat{\theta }$ -- ОМП.
\end{proposition}
\begin{proof}
    По неравенству Крамера-Рао $\displaystyle \hat{\theta }( X) -\theta =\dfrac{1}{I_{X}( \theta )}\left(\dfrac{\partial }{\partial \theta }\ln f_{\theta }( X)\right)$. Тогда, если $\displaystyle \theta < \hat{\theta }( X)$, то $\displaystyle \dfrac{\partial }{\partial \theta }\ln f_{\theta }( X)  >0\Rightarrow f_{\theta }( X)$ возрастает. Если же $\displaystyle \theta  >\hat{\theta }( X)$, то $\displaystyle f$ убывает. Следовательно, максимум функции правдоподобия достигается в точке $\displaystyle \hat{\theta }( X)$.
\end{proof}